\section{Introdução}

\paragraph{}
Nesta atividade foi realizado um processo de reconhecimento e exploração de
vetores de escalação de privilégios em um sistema Linux previamente
comprometido. O objetivo principal consistiu em identificar configurações
inseguras, permissões incorretas e mecanismos mal configurados que permitissem
a obtenção de privilégios administrativos no sistema.

\paragraph{}
Inicialmente foram utilizadas ferramentas automatizadas de enumeração para
coletar informações sobre o ambiente alvo, substituindo parte do processo de
reconhecimento manual realizado em atividades anteriores. A partir das
informações obtidas, foram analisadas diferentes oportunidades de exploração
relacionadas a arquivos críticos do sistema, permissões de execução com
\texttt{sudo} e acesso ao serviço Docker.

\paragraph{}
Durante os testes foram aplicadas técnicas clássicas de pós-exploração e
escalação de privilégios em ambientes Linux, demonstrando como pequenas falhas
de configuração podem resultar no comprometimento completo do sistema.
